
\documentclass[a4paper,12pt]{letter}
\usepackage[utf8]{inputenc}
\usepackage[english,slovene]{babel}
%\usepackage{blindtext}
%\usepackage{pdflscape}
\usepackage[tikz]{bclogo}
\usepackage{qrcode,marvosym}
\usepackage{pgf,tikz,pgfplots}
\pgfplotsset{compat=1.14}
\usepackage{mathrsfs}
\usetikzlibrary{arrows}
\usepackage{graphicx}
\newcommand{\ds}{\displaystyle}
\usepackage{fancybox}
\usepackage{amsfonts}
\usepackage{amsmath}
%\usepackage{color}
\usepackage{wallpaper}
\usepackage{hyperref}
\usepackage[cp1250]{inputenc}
\usepackage{array}%%%%%%%%%%%%%%%%%%%%%%%%%%%%%%%%%%%%%%%%%%%%
\newcolumntype{M}[1]{>{\centering\arraybackslash}m{#1}}%%%%
\newcolumntype{N}{@{}m{0pt}@{}}%%%%%%%%%%%%%%%%%%%%%%%%%%%%
   \usepackage{fancyhdr,enumerate}

 \newcommand{\ora}{\overrightarrow}
 \renewcommand{\headrulewidth}{3pt}
 \renewcommand{\headrule}{\hbox to\headwidth{%
   \color{gray}\leaders\hrule height \headrulewidth\hfill}}
    \renewcommand{\footrulewidth}{2pt}
    \renewcommand{\footrule}{\hbox to\headwidth{%
  \color{brown}\leaders\hrule height \footrulewidth\hfill}}
\usepackage{gensymb}
\usepackage{amssymb} 

  \usepackage{fancyhdr}
  \textwidth 20 true cm
  \oddsidemargin -2 true cm
  \evensidemargin -2 true cm
  \topmargin -3.5 true cm
  \textheight 27.5 true cm
  \linespread{1.2}
    \renewcommand{\headrulewidth}{0.3pt}
   \renewcommand{\footrulewidth}{1.9pt}
  \definecolor{qqwuqq}{rgb}{1,0,0 }
  \definecolor{qqqqff}{rgb}{0,0,1}
  \definecolor{qqffqq}{rgb}{0,1,0}
  \definecolor{ffqqqq}{rgb}{1,0,0}
\definecolor{zzuxux}{rgb}{0,0,0}%{0.92,0.03,0.03}
\usepackage{mdframed}
\usepackage{multicol}
 \definecolor{bgblue}{RGB}{0,0,253}
\usepackage{xcolor}
\usepackage{eurosym}
%%%%%%%%%%%%%%%%%%%%%%%%%%%%%%%%%%%%%%%%%%%%%%%%%%%%%%%%
\newcounter{tocke}
\setcounter{tocke}{0}
\usepackage{tikz-3dplot}
%%%%%%%%%%%%%%%%%%%%%%%%%%%%%%%%%%%%%%%%%%%%%%%%%%%%%%%%%%
\mdfdefinestyle{theoremstyle}{%
linecolor=brown!40,linewidth=1pt,%
frametitlerule=true,%
frametitlebackgroundcolor=brown!50,
innertopmargin=\topskip,
}

\mdfdefinestyle{theoremstyle1}{%
linecolor=brown,middlelinewidth=1pt,%
frametitlerule=true,%
apptotikzsetting={\tikzset{mdfframetitlebackground/.append style={%
shade,left color=black!40, right color=gray!10},
mdfbackground/.append style={
top color=white!100!white,
bottom color=black!5
},
}}, 
frametitlerulecolor=gray,
frametitlerulewidth=2pt,
innertopmargin=\topskip,
innerbottommargin=\topskip,
}
\mdtheorem[style=theoremstyle1]{naloga}{Naloga}
\mdtheorem[style=theoremstyle]{resitev}{Analiza Naloge}
\mdtheorem[style=theoremstyle]{kriterij}{\v{S}tevilo dose\v{z}enih to\v{c}k na testu}
%%%%%%%%%%%%%%%%%%%%%%%%%%%%%%%%%%%%%%%%%%%%%%%%%%%%%%%%%%
\usepackage[table]{xcolor}
\usepackage{venndiagram}
\usepackage[printwatermark]{xwatermark}
\usepackage{xcolor}
\newwatermark[allpages,color=brown!2,angle=45,scale=5,
xpos=-5,ypos=0]{matej.info}
%%%%%%%%%%%%%%%%%%%%%%%%%%%%%%%%%%%%%%%%%%%%%%%%%%%%%%%%%%
%%%%%%%%%%%%%%%%%%%%%%%%%%%%%%%%%%%%%%%%%%%%%%%%%%%%%%%%%%

%
\begin{document}
\pagestyle{fancy}

\begin{minipage}{20cm}
\begin{bclogo}[couleur=brown!40,ombre=true,logo=\bcplume,noborder=true]
{\Large Test 1.3: \large Polinomi in racionalne funkcije \hfill G-3}
\end{bclogo}
\end{minipage}

\vspace{0.3cm}

\begin{minipage}{20cm}
\begin{bclogo}[couleur=brown!40,ombre=true,logo=\bcplume,noborder=true]
{Ime in priimek:} \rule{15cm}{0.2mm}
\end{bclogo}
\end{minipage}

%%%%%%%%%%%%%%%%%%%%%%% 1. STRAN

\def\tockeA{3+2+3}
\begin{naloga}[\hfill $\tockeA$] \addtocounter{tocke}{\tockeA}
Podan je graf polinoma tretje stopnje.

a) Določi ničle in njihovo večkratnost.

b) Določi $p(0)$ in $p(2)$.

c) Zapiši predpis polinoma.


\definecolor{qqwuqq}{rgb}{0.,0.39215686274509803,0.}
\begin{tikzpicture}[line cap=round,line join=round,>=triangle 45,x=1.0cm,y=1.0cm]
\begin{axis}[
x=1.0cm,y=1.0cm,
axis lines=middle,
ymajorgrids=true,
xmajorgrids=true,
xmin=-4.6945454545454535,
xmax=3.814545454545456,
ymin=-5.3199999999999985,
ymax=2.407272727272729,
xtick={-4.0,-3.0,...,3.0},
ytick={-5.0,-4.0,...,2.0},]
\clip(-4.6945454545454535,-5.32) rectangle (3.814545454545456,2.407272727272729);
\draw[line width=1.pt,color=qqwuqq,smooth,samples=100,domain=-4.6945454545454535:3.814545454545456] plot(-\x,{((\x)-1)*((\x)+2)^(2)});
\begin{scriptsize}
\draw[color=qqwuqq] (-3.385454545454544,-8.783636363636361) node {$f$};
\end{scriptsize}
\end{axis}
\end{tikzpicture}
\end{naloga}

\vspace{8cm}
\newpage

\def\tockeB{5}
\begin{naloga}[\hfill $\tockeB$] \addtocounter{tocke}{\tockeB}
Izračunaj količnik in ostanek pri deljenju:
\[
p(x)=2x^3-3x^2+4x-5, \quad q(x)=x^2-x-2.
\]
\end{naloga}
\vfill


\vfill
\mdtheorem[style=theoremstyle1]{kriterij}{Kriterij ocenjevanja}
\begin{kriterij*}[\hfill \v{s}tevilo vseh to\v{c}k na testu: $\arabic{tocke}$]
 $$
 \begin{array}{||c|c|c|c|c|c||
>{}c|c|}
 \hline\hline
 \textup{ocena}& 1&2&3&4&5&\textrm{uspe\v{s}nost v } \%& \textrm{\bf OCENA}\\
  \hline\hline
 \% & [0,45) &[45,60)&[60,75
 ) &[75,90)    &    [90,100]& \mbox{\phantom{ \huge{$\frac{5}{445}$
 15}}
 } 
 & \\
 \hline
 \end{array} \hspace*{0.5cm}\vspace*{-0.3cm} \qrcode[hyperlink,height=0.8in]{http://www.matej.info}
$$
\vspace*{0.1cm}
\end{kriterij*}
\newpage

%%%%%%%%%%%%%%%%%%%%%%% 2. STRAN

\def\tockeC{5}
\begin{naloga}[\hfill $\tockeC$] \addtocounter{tocke}{\tockeC}
Določi vse ničle polinoma:
\[
p(x)=x^4-2x^3-7x^2+8x+12.
\]
Najprej zapiši vse možne celoštevilske kandidate.
\end{naloga}

\vspace{8cm}

\def\tockeD{3+3}
\begin{naloga}[\hfill $\tockeD$] \addtocounter{tocke}{\tockeD}
S Hornerjevim algoritmom izračunaj:

a) $p(3)$ za $p(x)=x^3-2x^2+4x-1$

b) Ali je $x^3+3x^2-4$ deljiv z $x-1$?
\end{naloga}

\newpage

%%%%%%%%%%%%%%%%%%%%%%% 3. STRAN

\def\tockeE{4}
\begin{naloga}[\hfill $\tockeE$] \addtocounter{tocke}{\tockeE}
Reši enačbo:
\[
x^3-5x^2+6x=0.
\]
\end{naloga}

\vspace{10cm}



%%%%%%%%%%%%%%%%%%%%%%% 4. STRAN

\def\tockeG{6}
\begin{naloga}[\hfill $\tockeG$] \addtocounter{tocke}{\tockeG}
Nariši graf funkcije:
\[
f(x)=\frac{x-2}{x-1}.
\]
Določi:
\begin{itemize}
\item ničlo
\item pol
\item asimptoti
\end{itemize}
\end{naloga}

\vfill


\newpage

\newpage
\section*{Rešitve}

\begin{resitev}
\textbf{1. naloga}

Podan je graf:
\[
f(x) = -x(x-1)(x+2)^2
\]

a) Ničle:
\[
x=0 \ (\text{enostavna}), \quad x=1 \ (\text{enostavna}), \quad x=-2 \ (\text{dvojna})
\]

b)
\[
p(0)=0,\quad p(2)= -2\cdot1\cdot 4 = -8
\]

c)
\[
p(x) = -x(x-1)(x+2)^2
\]
\end{resitev}

\begin{resitev}
\textbf{2. naloga}

Deljenje:
\[
\frac{2x^3-3x^2+4x-5}{x^2-x-2}
\]

Rezultat:
\[
k(x)=2x-1,\quad r(x)=2x-7
\]
\end{resitev}

\begin{resitev}
\textbf{3. naloga}

Možni kandidati:
\[
\pm1,\pm2,\pm3,\pm4,\pm6,\pm12
\]

Ničle:
\[
x=2,\quad x=-1,\quad x=3,\quad x=-2
\]

Torej:
\[
p(x)=(x-2)(x+1)(x-3)(x+2)
\]
\end{resitev}

\begin{resitev}
\textbf{4. naloga}

a)
\[
p(3)=27-18+12-1=20
\]

b)
\[
p(1)=1+3-4=0
\]

Polinom je deljiv z $x-1$.
\end{resitev}

\begin{resitev}
\textbf{5. naloga}

\[
x^3-5x^2+6x = x(x^2-5x+6)
\]

\[
= x(x-2)(x-3)
\]

Rešitve:
\[
x=0,2,3
\]
\end{resitev}

\begin{resitev}
\textbf{6. naloga}

\[
f(x)=\frac{x-2}{x-1}
\]

Ničla:
\[
x=2
\]

Pol:
\[
x=1
\]

Asimptoti:
\[
x=1 \ (\text{navpična}), \quad y=1 \ (\text{vodoravna})
\]
\end{resitev}

\end{document}